\subsection{Glossary}
\label{app:glossary}

\begin{description}
\item[Arithmetic intensity (AI)] FLOPs performed per byte fetched from
DRAM; the x-axis of the roofline. Low AI = memory-bound.
\item[Attribution] Assigning each memory access to the named data structure
(allocation) it falls within.
\item[Bandwidth] Data volume movable per unit time (GB/s); contrast
\emph{latency}.
\item[Bank (DRAM)] A subdivision of a DRAM device; near-bank PiM places
compute at each bank to harvest internal bandwidth.
\item[Bundle adjustment (BA)] Joint nonlinear least-squares refinement of
camera poses and 3-D landmarks; cuVSLAM's numerical core.
\item[Cache (L1/L2)] Automatic on-chip copies of recently used data; L1 per
SM, L2 chip-wide (25\,MB on the primary instrument).
\item[Coalescing] Hardware merging of a warp's 32 lane addresses into few
32-byte sectors; its failure mode (scatter) wastes up to 8$\times$
bandwidth.
\item[Cold-persistent] Data with session lifetime touched episodically
(the keyframe database).
\item[Compulsory miss] A first-touch cache miss no cache size can avoid;
streaming data consists almost entirely of these.
\item[CoV] Coefficient of variation (standard deviation $\div$ mean);
run-to-run noise measure (0.14\% at locked clocks here).
\item[DRAM speed-of-light (DRAM SoL)] Achieved DRAM bandwidth as a
percentage of the measured peak; $\geq$50\% is treated as saturation.
\item[Footprint] The number of distinct 32-byte sectors a launch touches,
i.e.\ its exact working-set size.
\item[G0--G7] The GPU bottleneck taxonomy (Table~\ref{tab:taxonomy}).
\item[Ground truth] The externally measured true trajectory against which
SLAM output is scored.
\item[Hot-persistent] Resident data reused every frame (the BA linear
system).
\item[Jaccard overlap] $|A \cap B| / |A \cup B|$ between consecutive
launches' footprints; measures working-set migration.
\item[Kernel] One GPU function launched over a grid of threads.
\item[Keyframe] A frame promoted into the long-term map; its descriptors
enter the loop-closure database.
\item[Lane] One of the 32 SIMT execution slots of a warp.
\item[LFMR] Last-to-first miss ratio; on a two-level GPU hierarchy,
$1 - \text{L2 hit rate}$. Near 1 = caches not helping.
\item[Local memory] Per-thread DRAM-backed space used for register spills;
coalesced by construction; \emph{not} program data.
\item[Loop closure] Recognition of a previously visited place, triggering
trajectory-wide drift correction via a database scan.
\item[Memory wall] The compounding gap between compute throughput and
memory speed; the thesis's founding observation.
\item[MPKI] DRAM sector fetches per thousand warp instructions.
\item[Occupancy] Resident warps as a fraction of the SM maximum; the GPU's
latency-hiding budget.
\item[Persistence] The lifetime/touch-pattern axis of a data structure:
streaming, hot-persistent, or cold-persistent.
\item[PiM / NDP] Processing-in-memory / near-data processing; compute
placed at or near DRAM.
\item[Reuse distance] Distinct sectors touched between successive accesses
to the same sector; its CDF predicts hit rate for any cache size.
\item[Roofline] $\min(\text{peak compute}, \text{AI} \times \text{peak
bandwidth})$ performance bound; placement under the sloped segment =
memory-bound.
\item[Sector] The 32-byte granule of GPU DRAM/L2 access; four per 128-byte
line.
\item[Sectors per request] Average sectors fetched per warp memory
request; $\sim$4 coalesced, $\rightarrow$32 scattered.
\item[SM] Streaming multiprocessor; the GPU's core unit (22 on the primary
instrument).
\item[Spill (register)] Compiler-generated traffic to local memory when
per-thread state exceeds the register budget.
\item[Stall reason] Hardware-recorded cause of a warp's failure to issue
(memory scoreboard, dependency wait, barrier, \ldots).
\item[Streaming] Per-frame data consumed once; compulsory-miss dominated.
\item[Substrate] A candidate execution home: GPU, CPU, DRAM-PiM,
scatter-PiM, near-sensor, in/near-storage.
\item[TaggedAllocator] This project's allocation-journaling
instrumentation: pointer, size, call stack, and data-structure tag for
every GPU allocation.
\item[Warp] The fixed 32-thread SIMT execution bundle.
\end{description}

\subsection{Artifact and Reproducibility Description}
\label{app:artifacts}

All software, configurations, and derived data of this thesis are archived
in a single repository (\code{github.com/mij001/cuvslam-stack}, MIT
license). The artifact is organized so that \emph{every table in this
document regenerates from committed CSV data using only the Python standard
library no GPU, no datasets, no vendor tools required}; capture (which
does require the hardware) is separated from analysis by construction.

\subsubsection{Layout}

\begin{description}
\item[\code{configs/}] The configuration regime: a small set of human-owned
base TOML files plus a deterministic mutation pass that materializes the
full experiment matrices (141 accuracy configurations, 192
profiling-coverage variants). One TOML fully describes a run: dataset,
camera rig, pipeline mode, feature toggles, and the ground-truth evaluation
block.
\item[\code{profiling/harness/}] The capture entry point: one command runs
any configuration under nsys, ncu (curated metric set, bounded windows), or
NVBit (windowed, kernel-filtered address traces), writing versioned result
directories with full provenance (clock locks, driver/tool versions,
command lines). A workload-\emph{adapter} interface makes the harness
generic: cuVSLAM is one adapter; a command adapter runs arbitrary GPU
programs through the identical pipeline.
\item[\code{profiling/analysis/}] The stdlib-only analysis library:
screening, classification (the decision tree of \S\ref{sec:met-tree}, with
thresholds defined once and imported live into documentation), roofline,
locality (reuse-distance CDFs with memory-space filtering), attribution
(the journal-trace join), campaign synthesis, substrate verdicts, residual
ranking, and the DAMOV-parity checks. The GPU-free test suite fabricates
known inputs for every module.
\item[\code{profiling/calibration/}] The known-truth archetype suite
(\S\ref{sec:char-calibration}): eight kernels, one per class, with the
blind-classification driver.
\item[\code{profiling/validation/}] The clock-domain intervention scripts
(\S\ref{sec:met-parity}) and the profiler-neutrality drivers.
\item[\code{patches/}] The instrumentation as reviewable diffs: the
TaggedAllocator + NVTX enablement patch against the cuVSLAM source tree,
and the NVBit \code{mem\_trace} extensions (launch-window and kernel-name
gating; the allocation-event sidecar).
\item[\code{reports/}, \code{profiling/reports/}] Dated, immutable result
directories: the 27-sequence characterization campaign, the locality study
and its space-filtered re-derivation (both the superseded and corrected
tables are retained, with the correction marked in place), the attribution
campaign, the 141-run accuracy matrix, the 192-variant neutrality campaign,
the substrate synthesis, the energy and host-I/O measurements, and the
DAMOV-validation artifacts (cross-device agreement, hierarchical-clustering
agreement).
\item[\code{dashboard/}] The evidence explorer: an interactive application
presenting every finding with its computed number, the per-run reasoning
chain (time share $\rightarrow$ metrics vs.\ thresholds $\rightarrow$ the
decision-tree rule that fired $\rightarrow$ verdict), and a methodology tab
whose formulas and thresholds are stamped from the analysis source at build
time so interface and code cannot disagree.
\end{description}

\subsubsection{Environment}

Primary testing hardware: NVIDIA RTX~2000 Ada (\code{sm\_89}), driver
575.64.05, CUDA 12.9.1, Nsight Compute 2025.2.1.3, Nsight Systems 2025.3.2,
and NVBit 1.8. Maintaining this exact combination of software and driver
versions is critical, since the binary instrumentation specifies the driver
version, which, in its turn, determines the profiler versions. To ensure
compatibility, an environment doctor checks the target system and suggests
a one-line fix for any incompatibilities detected. GPU clocks are locked
(1620/7001\,MHz) for every measurement; ceilings are re-measured per
instrument.

\subsubsection{Reproduction paths}

\textbf{Analysis-only (no GPU):} clone; run the test suite; regenerate any
table in this thesis from the committed CSVs via the corresponding
\code{analysis} module. \textbf{Full capture:} configure a target with the
environment doctor; run the regime's campaign drivers (all are resumable
and idempotent re-running skips completed steps); the coverage audit
then proves the captures complete before analysis will consume them.
\textbf{Instrumented builds:} apply the committed patch to the cuVSLAM
source tree and build with the pinned container recipe; the baseline and
instrumented binaries are kept side by side, and the QoR check of
\S\ref{sec:val-neutrality} re-verifies output equivalence on any rebuild.

\subsection{Extended Result Tables}
\label{app:tables}

The tables below are excerpts of committed artifacts, reproduced for
convenience; the full machine-readable versions live in the repository
(Appendix~\ref{app:artifacts}).

\subsubsection{Cross-sequence attribution consistency (excerpt)}

Selected rows of the 49-kernel attribution-consistency table: each kernel's
modal top global-space tag, the agreement across every sequence it appears
in, and its median memory-space split (shared/spill/global).

\begin{table}[h]
\centering
\caption{Attribution consistency across 27 sequences (excerpt).}
\label{tab:attr-consistency}
\small
\begin{tabular}{@{}llcccc@{}}
\toprule
kernel & modal global tag & agree & shared & spill & global \\
\midrule
\code{cast\_image\_kernel} & \code{images\_raw} & 100\% & 0\% & 0\% & 100\% \\
\code{cast\_image\_kernel\_rgb} & \code{pyramid\_levels} & 100\% & 0\% & 0\% & 100\% \\
\code{gaussian\_scaling\_kernel} & \code{images\_raw} & 100\% & 0\% & 0\% & 100\% \\
\code{conv\_grad\_x\_kernel} & \code{pyramid\_levels} & 100\% & 92\% & 0\% & 8\% \\
\code{conv\_grad\_y\_kernel} & \code{pyramid\_levels} & 100\% & 96\% & 0\% & 4\% \\
\code{gftt\_values\_kernel} & \code{feature\_tracks} & 100\% & 98\% & 0\% & 2\% \\
\code{non\_max\_suppression} & \code{feature\_tracks} & 100\% & 90\% & 0\% & 10\% \\
\code{accumulateGFTT\_kernel} & \code{feat.\ sel.\ scratch} & 100\% & 89\% & 0\% & 11\% \\
\code{sba::build\_full\_system\_1} & \code{ba\_linear\_system} & 100\% & 0\% & 0\% & 100\% \\
\code{sba::calc\_jacobians} & \code{ba\_linear\_system} & 100\% & 4\% & 0\% & 96\% \\
\code{sba::reduced\_system\_st.\,2} & \code{ba\_linear\_system} & 100\% & 4\% & 0\% & 96\% \\
\stBuild{} & \code{keyframe\_descriptors} & 100\% & 0\% & 0\% & 100\% \\
\stTrack{} & \code{keyframe\_descriptors} & 100\% & 0\% & 93\% & 7\% \\
\code{cub::DeviceMergeSort*} & \code{feature\_tracks} & 100\% & 83--97\% & 0\% & 3--17\% \\
\bottomrule
\end{tabular}
\end{table}

\subsubsection{Instrumented vs.\ baseline trajectories (QoR excerpt)}

\begin{table}[h]
\centering
\caption{Quality-of-results check: instrumented (TaggedAllocator) build vs.\
baseline build, same configurations (APE, m).}
\label{tab:qor}
\begin{tabular}{@{}lccc@{}}
\toprule
configuration & baseline & instrumented & $\Delta$ \\
\midrule
EuRoC MH\_01 inertial SLAM & 0.689 & 0.689 & 0.0000 \\
EuRoC V1\_01 stereo SLAM & 0.037 & 0.037 & 0.0000 \\
EuRoC V2\_02 stereo odometry & 0.177 & 0.177 & 0.0000 \\
KITTI 00 stereo odometry & 6.841 & 6.605 & $-$0.236\textsuperscript{a} \\
KITTI 06 stereo SLAM & 2.312 & 2.277 & $-$0.035\textsuperscript{a} \\
TUM long\_office RGB-D SLAM & 0.018 & 0.021 & $+$0.003 \\
\bottomrule
\end{tabular}

\smallskip
\raggedright\footnotesize \textsuperscript{a}\,Within the sequence's own
plain-rerun scatter (\S\ref{sec:val-nondeterminism}); kitti00 is the
canonical bimodal loop-closure case.
\end{table}

\subsubsection{KITTI per-sequence odometry accuracy}

\begin{table}[h]
\centering
\caption{KITTI stereo odometry, per sequence: absolute error and
segment-relative drift (100--800\,m KITTI protocol).}
\label{tab:kitti}
\begin{tabular}{@{}lccc@{}}
\toprule
seq. & APE (m) & avgRTE (\%) & avgRE ($^\circ$) \\
\midrule
00 & 6.84 & 0.892 & 0.938 \\
01 & 11.59 & 1.899 & 0.386 \\
02 & 4.55 & 0.975 & 0.707 \\
03 & 3.94 & 2.265 & 0.988 \\
04 & 2.16 & 1.967 & 0.090 \\
05 & 3.31 & 0.974 & 0.723 \\
06 & 2.88 & 1.248 & 0.757 \\
07 & 1.44 & 0.982 & 1.131 \\
08 & 4.57 & 1.221 & 0.959 \\
09 & 3.53 & 1.238 & 0.544 \\
10 & 2.17 & 1.010 & 0.661 \\
\bottomrule
\end{tabular}

\smallskip
\raggedright\footnotesize At the 500\,m segment length the corpus-level
translational drift is 0.82\%, against the 0.85\% published KITTI-protocol
figure (\S\ref{sec:val-accuracy}). Sequence 01 (highway, sparse features)
dominates the absolute-APE average; sequences 01 and 04 close no loops and
serve as the loop-free control in the attribution campaign.
\end{table}

\subsubsection{Measured stage map (NVTX, excerpt)}

\begin{table}[h]
\centering
\caption{Kernel-to-stage mapping, measured from cuVSLAM's own NVTX ranges.}
\label{tab:nvtx}
\small
\begin{tabular}{@{}lp{0.55\textwidth}@{}}
\toprule
NVTX stage & kernels \\
\midrule
front-end, per frame & \code{cast\_image\_kernel\_rgb},
\code{gaussian\_scaling}, \code{conv\_grad\_x/y}, \code{lk\_track} \\
keyframe feature refresh & \code{gftt\_values}, \code{accumulateGFTT},
\code{downsample\_gftt}, \code{non\_max\_suppression},
\code{select\_features}, CUB keypoint sort \\
visual-odometry solver & \code{matcher::*}, \code{cast\_depth},
\code{sba::reduced\_system\_*}, cuSOLVER helpers \\
windowed bundle adjustment & \code{sba::build\_full\_system\_*},
\code{sba::calc\_jacobians}, \code{sba::evaluate\_cost} \\
keyframe ingest & \stBuild{} \\
loop closure \& optimization & \stTrack{} \\
\bottomrule
\end{tabular}
\end{table}
